\begin{figure}[htbp]
			\centering
			\resizebox{\textwidth}{!}{%
				\begin{tikzpicture}[
					scale=0.95,
					>=latex,
					trave/.style={line width=1.2pt},
					carico/.style={black, line width=0.6pt},
					forza/.style={black, line width=1pt},
					diagT/.style={black, line width=0.8pt},
					diagM/.style={black, line width=0.8pt},
					hatch/.style={gray!50, line width=0.3pt},
					quota/.style={black, line width=0.4pt},
					fibre/.style={gray!70, line width=0.3pt},
					]
					
					% ── Parametri dei vincoli (convenzione C06/C07) ──
					\def\tw{0.15}    \def\thR{0.28}   \def\rn{0.06}
					\def\rw{0.05}    \def\bw{0.09}    \def\gw{0.22}
					\def\hh{0.10}
					
					% ── Geometria del problema ──
					\def\Lell{4}                            % lunghezza trave ℓ
					\pgfmathsetmacro{\xA}{0}
					\pgfmathsetmacro{\xZero}{\Lell/sqrt(3)} % z₀ = ℓ/√3 ≈ 2.309 (punto T=0)
					\def\yqi{-0.80}                         % y delle quote sotto la trave
					\def\ylab{-3.30}                        % y delle etichette di pannello
					
					% ── Parametri del carico distribuito ──
					\def\pgap{0.05}                         % gap fra trave e base del triangolo
					\def\php{0.40}                          % altezza massima (in B)
					\pgfmathsetmacro{\pBase}{\pgap}         % base del triangolo: y = 0.05
					\pgfmathsetmacro{\pTop}{\pgap+\php}     % vertice destro: y = 0.45
					
					% =========================================================
					% PANNELLO (a) — Schema statico
					% =========================================================
					\begin{scope}[xshift=0cm]
						
						% ── Trave ──
						\draw[trave] (0,0) -- (\Lell,0);
						
						% ── Fibra inferiore (tratteggio sotto la trave) ──
						\draw[fibre, dashed] (0, -0.1) -- (\Lell, -0.1);
						
						% ── Cerniera in A ──
						\draw[thick] (0,0) -- (-\tw,-\thR) -- (\tw,-\thR) -- cycle;
						\draw[thick] (-\gw,-\thR) -- (\gw,-\thR);
						\foreach \x in {-0.16,-0.08,0,0.08,0.16}
						\draw[thin] (\x,-\thR) -- ++(-0.07,-\hh);
						
						% ── Carrello verticale in B ──
						\draw[thick] (\Lell,0) -- (\Lell-\tw,-\thR) -- (\Lell+\tw,-\thR) -- cycle;
						\draw[thick] (\Lell-\bw,-\thR-\rw-0.01) circle (\rw);
						\draw[thick] (\Lell+\bw,-\thR-\rw-0.01) circle (\rw);
						\draw[thick] (\Lell-\gw,-\thR-2*\rw-0.04) -- (\Lell+\gw,-\thR-2*\rw-0.04);
						\foreach \x in {-0.16,-0.08,0,0.08,0.16}
						\draw[thin] (\Lell+\x,-\thR-2*\rw-0.04) -- ++(-0.07,-\hh);
						
						% ── Anellini bianchi sui perni ──
						\draw[thick, fill=white] (0,0) circle (\rn);
						\draw[thick, fill=white] (\Lell,0) circle (\rn);
						
						% ── Etichette A, B ──
						\node[above left=-2pt,  font=\small] at (0,0)      {$A$};
						\node[above right=-2pt, font=\small] at (\Lell,0)  {$B$};
						
						% ── Carico distribuito triangolare p(z) = p̄·z/ℓ ──
						% Triangolo: base orizzontale sopra la trave, ipotenusa da A a B
						\draw[carico] (0,\pBase)     -- (\Lell,\pBase);      % base
						\draw[carico] (0,\pBase)     -- (\Lell,\pTop);       % ipotenusa
						\draw[carico] (\Lell,\pBase) -- (\Lell,\pTop);       % verticale destra
						% Frecce: lunghezza linearmente crescente da 0 in A a piena in B
						% Quota dell'ipotenusa a x: y = pBase + (pTop-pBase)·(x/ℓ)
						\foreach \x in {1.0, 1.5, ..., 4.0} {\pgfmathsetmacro{\ytop}{\pBase + (\pTop-\pBase)*(\x/\Lell)}
							\draw[->, carico, >={Latex[length=2.5pt, width=2.0pt]}]
							(\x, \ytop) -- (\x, \pBase);
						}
						
						% Etichetta p̄ in alto a destra, vicino al vertice del triangolo
						\node[above, font=\small] at (\Lell+0.05, \pTop) {$\bar{p}$};
						
						% ── Quota ℓ sotto la trave ──
						\draw[<->, quota] (0,\yqi) -- (\Lell,\yqi);
						\draw[quota] (0,\yqi+0.07)     -- (0,\yqi-0.07);
						\draw[quota] (\Lell,\yqi+0.07) -- (\Lell,\yqi-0.07);
						\node[above, font=\small] at (0.5*\Lell, \yqi-0.04) {$\ell$};
						
						% Etichetta pannello (a)
						\node[font=\small] at (0.5*\Lell, 0.4*\ylab) {(a)};
						
					\end{scope}
					
					% =========================================================
					% PANNELLO (b) — Diagramma di T(z)
					% =========================================================
					\begin{scope}[xshift=5.5cm]
						
						% ── Trave (clone) ──
						\draw[trave] (0,0) -- (\Lell,0);
						
						% ── Vincoli (clone) ──
						\draw[thick] (0,0) -- (-\tw,-\thR) -- (\tw,-\thR) -- cycle;
						\draw[thick] (-\gw,-\thR) -- (\gw,-\thR);
						\foreach \x in {-0.16,-0.08,0,0.08,0.16}
						\draw[thin] (\x,-\thR) -- ++(-0.07,-\hh);
						\draw[thick] (\Lell,0) -- (\Lell-\tw,-\thR) -- (\Lell+\tw,-\thR) -- cycle;
						\draw[thick] (\Lell-\bw,-\thR-\rw-0.01) circle (\rw);
						\draw[thick] (\Lell+\bw,-\thR-\rw-0.01) circle (\rw);
						\draw[thick] (\Lell-\gw,-\thR-2*\rw-0.04) -- (\Lell+\gw,-\thR-2*\rw-0.04);
						\foreach \x in {-0.16,-0.08,0,0.08,0.16}
						\draw[thin] (\Lell+\x,-\thR-2*\rw-0.04) -- ++(-0.07,-\hh);
						\draw[thick, fill=white] (0,0) circle (\rn);
						\draw[thick, fill=white] (\Lell,0) circle (\rn);
						
						% ── Etichette A, B ──
						\node[above left=-2pt,  font=\small] at (0,0)      {$A$};
						\node[above right=-2pt, font=\small] at (\Lell,0)  {$B$};
						
						% ── Tacca verticale in z₀ = ℓ/√3 ──
						\draw[forza] (\xZero, -0.10) -- (\xZero, 0.10);
						
						% ── Diagramma T(z) ──
						\def\sT{0.75}
						\pgfmathsetmacro{\TA}{ (\Lell/6)*\sT}                % T(A) = +p̄ℓ/6 → +0.30
						\pgfmathsetmacro{\TB}{-(\Lell/3)*\sT}                % T(B) = -p̄ℓ/3 → -0.60
						
						% Hatching verticale: dall'asse al diagramma — formula T(x) normalizzata
						\foreach \x in {0.20, 0.40, ..., 3.80} {
							\pgfmathsetmacro{\h}{\sT*(\Lell/6 - \x*\x/(2*\Lell))}
							\draw[hatch] (\x, 0) -- (\x, \h);
						}
						
						% Bordo del diagramma — parabola campionata a passo 0.20
						\draw[diagT]
						(0, 0) -- (0, \TA)
						\foreach \x in {0.20, 0.40, ..., 3.80} {
							-- ({\x}, {\sT*(\Lell/6 - \x*\x/(2*\Lell))})
						}
						-- (\Lell, \TB) -- (\Lell, 0);
						
						% Segni cerchiati nelle due aree
						\node[font=\small] at ({0.4*\xZero}, {0.4*\TA})	{$\boldsymbol{\oplus}$};
						\node[font=\small] at ({0.55*(\xZero+\Lell)}, {0.4*\TB})	{$\boldsymbol{\ominus}$};
						
						% Valori notevoli (senza segno, come da convenzione)
						\node[above left,  font=\small] at (0, \TA-0.2)     {$\tfrac{\bar{p}\ell}{6}$};
						\node[right, font=\small] at (\Lell, \TB) {$\tfrac{\bar{p}\ell}{3}$};
						
						% ── Quota di z₀ sotto la trave (da A a z₀) ──
						\draw[<->, quota] (0,\yqi)      -- (\xZero,\yqi);
						\draw[quota] (0,\yqi+0.07)      -- (0,\yqi-0.07);
						\draw[quota] (\xZero,\yqi+0.07) -- (\xZero,\yqi-0.07);
						\node[above, font=\footnotesize] at ({0.5*\xZero}, \yqi-0.04)	{$\ell/\sqrt{3}$};
						
						% Etichetta pannello (b)
						\node[font=\small] at (0.5*\Lell, 0.4*\ylab) {(b)};
						
					\end{scope}
					
					% =========================================================
					% PANNELLO (c) — Diagramma di M(z)
					% =========================================================
					\begin{scope}[xshift=11.0cm]
						
						% ── Trave (clone) ──
						\draw[trave] (0,0) -- (\Lell,0);
						
						% ── Vincoli (clone) ──
						\draw[thick] (0,0) -- (-\tw,-\thR) -- (\tw,-\thR) -- cycle;
						\draw[thick] (-\gw,-\thR) -- (\gw,-\thR);
						\foreach \x in {-0.16,-0.08,0,0.08,0.16}
						\draw[thin] (\x,-\thR) -- ++(-0.07,-\hh);
						\draw[thick] (\Lell,0) -- (\Lell-\tw,-\thR) -- (\Lell+\tw,-\thR) -- cycle;
						\draw[thick] (\Lell-\bw,-\thR-\rw-0.01) circle (\rw);
						\draw[thick] (\Lell+\bw,-\thR-\rw-0.01) circle (\rw);
						\draw[thick] (\Lell-\gw,-\thR-2*\rw-0.04) -- (\Lell+\gw,-\thR-2*\rw-0.04);
						\foreach \x in {-0.16,-0.08,0,0.08,0.16}
						\draw[thin] (\Lell+\x,-\thR-2*\rw-0.04) -- ++(-0.07,-\hh);
						\draw[thick, fill=white] (0,0) circle (\rn);
						\draw[thick, fill=white] (\Lell,0) circle (\rn);
						
						% ── Etichette A, B ──
						\node[above left=-2pt,  font=\small] at (0,0)      {$A$};
						\node[above right=-2pt, font=\small] at (\Lell,0)  {$B$};
						
						% ── Tacca verticale in z₀ = ℓ/√3 (punto di massimo) ──
						\draw[forza] (\xZero, -0.10) -- (\xZero, 0.10);
						
						% ── Diagramma M(z) ──
						\def\sM{0.80}
						
						% Hatching verticale: dall'asse al diagramma
						\foreach \x in {0.20, 0.40, ..., 3.80} {
							\pgfmathsetmacro{\h}{-\sM*(\Lell*\x/6 - \x*\x*\x/(6*\Lell))}
							\draw[hatch] (\x, 0) -- (\x, \h);
						}
						
						% Bordo del diagramma — cubica campionata a passo 0.20
						\draw[diagM]
						(0, 0)
						\foreach \x in {0.20, 0.40, ..., 3.80} {
							-- ({\x}, {-\sM*(\Lell*\x/6 - \x*\x*\x/(6*\Lell))})
						}
						-- (\Lell, 0);
						
						% Valore notevole — etichetta sotto il vertice della cubica in z₀
						\pgfmathsetmacro{\Mmax}{-\sM*\Lell*\Lell/(9*sqrt(3))}
						\node[above, font=\small] at (\xZero, \Mmax-0.04)	{$\tfrac{\bar{p}\ell^2}{9\sqrt{3}}$};
						
						% Etichetta pannello (c)
						\node[font=\small] at (0.5*\Lell, 0.4*\ylab) {(c)};
						
					\end{scope}
					
				\end{tikzpicture}
			}
			
			\label{fig:stat_trasv_iso_04}
		\end{figure}
